\begin{abstract}
% \ac{fca} enables the generation of concept hierarchies without requiring an expert ontology.
% This raises the question of whether the resulting hierarchy aligns with existing sources of knowledge. 
% To guide the construction of coherent concept hierarchies, prior work in constrained clustering introduces the use of hierarchical \ac{mlb} constraints.

    % 1. build a world: FCA and constrained clustering
    % 2. stress the problem: how do different knowledge sources align?
\ac{fca} enables the construction of concept hierarchies from formal contexts without requiring expert ontologies. 
Hierarchical expert knowledge can be encoded through \ac{mlb} constraints in constrained clustering.
However, it remains unclear how agreement between hierarchies derived from heterogeneous knowledge sources can be quantified systematically.

% first contribution
In this work, we prove that \acs{mlb} constraints induce a valid closure operator, providing a theoretical foundation for incorporating hierarchical expert knowledge into \ac{fca}. 
% second contribution
Building on this result, we introduce a method for comparing concept lattices derived from \acs{mlb} constraints and data-driven topic models. 
% case study
Using the BankSearch dataset, we compare concept lattices derived from expert-defined \acs{mlb} constraints and a \acl{lda}-based document-topic context. We analyze their agreement through \textit{consistency} and \textit{coherence} lattice intersections at different levels of coherence.
Our analysis reveals distinct structural differences between both hierarchies, highlighting the need for relaxed coherence measures when integrating heterogeneous knowledge sources.

% redundant
% In this work, we introduce a method for comparing and integrating knowledge derived from \acs{mlb} constraints with an additional hierarchical knowledge source. 
% To evaluate the benefits of our approach, we conduct experiments on the BankSearch dataset and explore the universe of a closure operator provided by the \acs{mlb} constraints.
% These constraints are derived from a manually annotated topic hierarchy and subsequently mapped onto the corresponding document corpus.
% As our supplementary hierarchical knowledge source, we employ a \acl{lda} topic model, which constructs a document-topic context through statistical inference over the same corpus.

% We mathematically prove that MLB constraints induce a valid closure operator, providing the theoretical foundation for integrating hierarchical expert knowledge into FCA. 
% We then analyze the structural differences between the concept lattices guided by the two knowledge sources. 
% The influence of both knowledge bases is assessed qualitatively through manual review and quantitatively by analyzing their intersection for different levels of required coherence.
% We find distinct structural variations between the two concept lattices, reflected in their \textit{consistency} and \textit{coherence} lattice intersection.
% Our findings suggest a significant divergence between expert-defined hierarchies and data-driven topics, highlighting the necessity of relaxed coherence measures for practical knowledge integration.

\keywords{\acl{fca} \and Constraints \and Topic Modeling.}
\end{abstract}